Es schmilzt Eis (Dichte $\rho_\text{E} = \rhoEisO$, FoTaBe S.~187) der Masse
\begin{align}
 m &= \mEisF = \rhoEis\cdot\AAr\cdot\hEis = \mEis
\end{align}

\emph{Arktis.} Das Eis schwimmt und verdrängt deshalb Meerwasser seiner eigenen Masse, also das Volumen $m/\rho_\text{M}$. Geschmolzen wird es zu Süsswasser mit dem Volumen $m/\rho_\text{W}$ ($\rho_\text{W} = \rhoWO$ beim Schmelzpunkt, FoTaBe S.~192). Weil Süsswasser weniger dicht ist als Meerwasser, braucht das Schmelzwasser etwas mehr Platz als das verdrängte Meerwasser. Verteilt auf die Fläche der Weltmeere steigt der Meeresspiegel um
\begin{align}
 \Delta h &= \frac{m}{A_\text{M}}\left(\frac{1}{\rho_\text{W}}-\frac{1}{\rho_\text{M}}\right) = \dhArF \\
   &= \frac{\mEis}{\AMe}\cdot\left(\frac{1}{\rhoW}-\frac{1}{\rhoMe}\right) \\
   &= \dhAr \approx \result{\dhArP{0}{2}}
\end{align}
Würde das Eis in Süsswasser schwimmen, änderte sich der Spiegel überhaupt nicht (wie beim Eiswürfel im Wasserglas).

\emph{Antarktis.} Eis auf dem Festland verdrängt kein Meerwasser. Sein ganzes Schmelzwasser fliesst zusätzlich ins Meer. Dieselbe Eismenge würde den Meeresspiegel um
\begin{align}
 \Delta h_\text{Land} &= \dhAnF = \frac{\mEis}{\rhoW\cdot\AMe} = \dhAn \approx \result{\dhAnP{0}{2}}
\end{align}
anheben, das ist rund \qAnP{0}{2}-mal mehr. Zudem liegt auf der Antarktis viel mehr Eis.
